\documentclass[a4paper]{article}

% Basic packages
\usepackage[fontset=fandol]{ctex}
\usepackage{amsmath, amssymb}
\usepackage{graphicx}
\usepackage[margin=2.5cm]{geometry}
\usepackage[most]{tcolorbox}
\usepackage{etoolbox}
\usepackage{listings}
\usepackage{hyperref}
\usepackage{booktabs}
\usepackage{subcaption}
\usepackage{float}
\usepackage{tikz}
\IfFileExists{pgfplots.sty}{
    \usepackage{pgfplots}
    \pgfplotsset{compat=1.18}
}{}

% ===== Highlight boxes =====
\newtcolorbox{knowledgebox}[1]{
    enhanced, colback=blue!5!white, colframe=blue!75!black, colbacktitle=blue!75!black,
    coltitle=white, fonttitle=\bfseries, title=#1, attach boxed title to top left={yshift=-2mm, xshift=2mm},
    boxrule=1pt, sharp corners
}

\newtcolorbox{importantbox}[1]{
    enhanced, colback=yellow!10!white, colframe=yellow!80!black, colbacktitle=yellow!80!black,
    coltitle=black, fonttitle=\bfseries, title=#1, sharp corners
}

\newtcolorbox{warningbox}[1]{
    enhanced, colback=red!5!white, colframe=red!75!black, colbacktitle=red!75!black,
    coltitle=white, fonttitle=\bfseries, title=#1, sharp corners
}

% ===== Code listings =====
\lstset{
    basicstyle=\ttfamily\small,
    keywordstyle=\color{blue},
    stringstyle=\color{red!60!black},
    commentstyle=\color{green!60!black},
    breaklines=true,
    frame=single,
    numbers=left,
    numberstyle=\tiny\color{gray},
    captionpos=b,
    extendedchars=false
}

% ===== Front-page metadata =====
% Fill these in when generating notes.
\newcommand{\notetitle}{课程笔记：[课程主题]}
\newcommand{\notesubtitle}{}
\newcommand{\noteauthors}{lecture-to-notes}
\newcommand{\notedate}{\today}
\newcommand{\videochannel}{[频道/作者]}
\newcommand{\videopublishdate}{[发布日期]}
\newcommand{\videoduration}{[视频时长]}
\newcommand{\videourl}{}
\newcommand{\videocoverpath}{}

\begin{document}

% ===== Title page =====
\begin{titlepage}
\centering
{\Large 课程笔记\par}
\vspace{1.2cm}
{\huge\bfseries \notetitle\par}
\ifdefempty{\notesubtitle}{}{
\vspace{0.3cm}
{\large \notesubtitle\par}
}
\vspace{0.5cm}
{\large \noteauthors\par}
\vspace{0.3cm}
{\large \notedate\par}
\vspace{1.0cm}

\ifdefempty{\videocoverpath}{
{\small 请在生成笔记时填入视频封面图的本地路径。\par}
}{
\includegraphics[width=0.82\textwidth,height=0.40\textheight,keepaspectratio]{\videocoverpath}\par
}

\vfill
\begin{tcolorbox}[width=0.9\textwidth, colback=black!2!white, colframe=black!60, sharp corners]
\textbf{视频作者/频道}：\videochannel\par
\textbf{发布日期}：\videopublishdate\par
\textbf{视频时长}：\videoduration\par
\textbf{视频链接}：
\ifdefempty{\videourl}{
未填写
}{
\href{\videourl}{\nolinkurl{\videourl}}
}
\end{tcolorbox}
\end{titlepage}

\tableofcontents
\newpage

%% --- 正文内容开始 --- %%

% INSTRUCTIONS FOR THE GENERATING AGENT:
%
% 1. Fill the metadata commands above (\notetitle, \videocoverpath, etc.)
% 2. Replace this block with the generated notes.
%
% FIGURE PATTERN (use smart-cropped frames):
%   \begin{figure}[H]
%   \centering
%   \includegraphics[width=\textwidth]{figures/fig_xxx.png}
%   \caption{描述\protect\footnotemark}
%   \end{figure}
%   \footnotetext{视频画面时间区间：00:12:31--00:12:46。}
%
% BOX RULES:
%   - importantbox: core concepts, definitions, key mechanisms
%   - knowledgebox: background, history, design tradeoffs, analogies
%   - warningbox:   common mistakes, hidden assumptions, pitfalls
%   - NO figures inside boxes
%   - No quota per section — use as many as the teaching signal demands
%
% SECTION STRUCTURE:
%   - Every major \section ends with \subsection{本章小结}
%   - Optionally add \subsection{拓展阅读} for external links
%   - Final section: \section{总结与延伸} with speaker summary + your synthesis
%
% WRITING RULES:
%   - Default language: Chinese
%   - Reconstruct teaching flow, don't mirror subtitle order
%   - Math: $$...$$ then symbol explanation list
%   - Code: lstlisting with caption
%   - No [cite] placeholders

%% --- 正文内容结束 --- %%

\end{document}
